\documentclass[12pt]{article}
\usepackage[margin=1in]{geometry}
\usepackage{amsmath,amsthm,amssymb,scrextend}
\setlength{\parindent}{0pt}
\begin{document}
\newcommand{\ds}{\displaystyle}

\textbf{(i)}

\begin{align*}
\frac{d}{dx}\left(\frac{P(x)}{Q(x)}e^x\right)
&=\left(\frac{P(x)}{Q(x)}\right)'e^x+\frac{P(x)}{Q(x)}e^x\\
&=\frac{P'(x)Q(x)-P(x)Q'(x)}{Q(x)^2}e^x+\frac{P(x)Q(x)}{Q(x)^2}e^x\\
&=\frac{e^x\left(Q(x)(P(x)+P'(x))-P(x)Q'(x)\right)}{Q(x)^2}\\
&=\frac{x^3-2}{(x+1)^2}e^x
\end{align*}


$$\therefore \bigl[Q(x)(P(x)+P'(x))-P(x)Q'(x)\bigr](x+1)^2=(x^3-2)Q(x)^2$$

When $\ds x=-1\implies0=-3Q(-1)^2\implies Q(-1)=0$

$$\therefore \boxed{x+1\text{ is a factor of }Q(x)}$$

Let $\ds\deg(P(x))=p,\ \deg(Q(x))=q \implies \deg(P(x)+P'(x))=p,\  \deg\bigl(Q(x)(P(x)+P'(x))\bigr)=p+q, \  \deg\bigl(P(x)Q'(x)\bigr)=p+q-1$

$$\deg\left(Q(x)(P(x)+P'(x))-P(x)Q'(x)\right)=p+q$$

$$\therefore \deg\left(\bigl[Q(x)(P(x)+P'(x))-P(x)Q'(x)\bigr](x+1)^2\right)=p+q+2$$

$$\deg\left((x^3-2)Q(x)^2\right)=2q+3$$

$$p+q+2=2q+3\implies p=q+1$$

$$\therefore \boxed{\deg(P(x))=\deg(Q(x))+1}$$

When $Q(x)=x+1$, $Q'(x)=1$

\begin{align*}
\bigl[Q(x)(P(x)+P'(x))-P(x)Q'(x)\bigr](x+1)^2&=(x^3-2)Q(x)^2\\
\bigl[(x+1)(P(x)+P'(x))-P(x)\bigr](x+1)^2&=(x^3-2)(x+1)^2\\
(x+1)(P(x)+P'(x))-P(x)&=x^3-2\\
xP(x)+(x+1)P'(x)&=x^3-2
\end{align*}

Let $P(x)=ax^2+bx+c$
\begin{align*}
xP(x)+(x+1)P'(x)
&=x(ax^2+bx+c)+(x+1)(2ax+b)\\
&=ax^3+(b+2a)x^2+(c+2a+b)x+b
\end{align*}

$$
\left\{
\begin{aligned}
a&=1\\
b+2a&=0\\
c+2a+b&=0\\
b&=-2
\end{aligned}
\right.
\quad\Longrightarrow\quad
\left\{
\begin{aligned}
a&=1\\
b&=-2\\
c&=0
\end{aligned}
\right.
$$

$$\therefore \boxed{P(x)=x^2-2x}$$

\bigskip\hrule\bigskip

\textbf{(ii)}

Assume such $P(x),Q(x)$ exist and $\text{gcd}(P,Q)=1$

\begin{align*}
\frac{d}{dx}\left(\frac{P(x)e^x}{Q(x)}\right)&=\frac1{x+1}e^x\\
\frac{e^x\left(Q(x)(P(x)+P'(x))-P(x)Q'(x)\right)}{Q(x)^2}&=\frac1{x+1}e^x
\end{align*}

$$\therefore \Bigl[Q(x)(P(x)+P'(x))-P(x)Q'(x)\Bigr](x+1)=Q(x)^2$$

When $\ds x=-1\implies0=Q(-1)^2\implies Q(-1)=0$

Let $Q(x)=(x+1)R(x)$

$$Q'(x)=R(x)+(x+1)R'(x)$$

$$
(x+1)R(x)(P(x)+P'(x))-P(x)(R(x)+(x+1)R'(x))=(x+1)R(x)^2
$$

When $\ds x=-1\implies -P(-1)R(-1)=0$

$\because \ds \text{gcd}(P,Q)=1$ and $x+1\mid Q(x) \quad \therefore \ds R(-1)=0 \implies (x+1)^2\mid Q(x)$

Let $\ds Q(x)=(x+1)^nS(x),\qquad n\geq2,\qquad S(-1)\neq0$

$$Q'(x)=n(x+1)^{n-1}S(x)+(x+1)^nS'(x)$$

$$\Bigl[Q(x)(P(x)+P'(x))-P(x)Q'(x)\Bigr](x+1)=Q(x)^2$$

$$
(x+1)^n\Bigl[(x+1)S(x)(P(x)+P'(x))-P(x)(nS(x)+(x+1)S'(x))\Bigr]=(x+1)^{2n}S(x)^2
$$

$$(x+1)S(x)(P(x)+P'(x))-P(x)(nS(x)+(x+1)S'(x))=(x+1)^nS(x)^2$$

When $x=-1\implies -nP(-1)S(-1)=0$

$$\because n\neq0,\qquad P(-1)\neq0,\qquad S(-1)\neq0$$

Contradiction

$$\therefore \boxed{\text{No such polynomials }P(x),Q(x)}$$

\bigskip
\hrule
\begin{center}
    This solution was submitted by \textbf{Jon}\\[4pt]
    Published and refined by \textbf{Jon}\\[8pt]
    \textit{Last updated: June 23, 2026}
\end{center}

\end{document}
